\subsection{Quantitative transfer from eigenvalue error to mode-root error}
\label{ssec:quantitative-mode-transfer}

This result concerns the genuine fixed-cutoff mode channel. It sharpens the
generic dual identity of Section~\ref{ssec:indexed-resolution-budget}
using the CCM divided-difference structure. The use of adjoint-weighted
Galerkin residuals is classical~\cite{IndexedBeckerRannacher}; the estimate
below exploits the two inverse-frequency powers in the even adjoint tail.
The finite transform is that of~\cite[Prop.~5.9]{CCM2025}. Qualitative
ground-state convergence under the hypotheses of~\cite[Thm.~6.1]{CvS2025}
does not itself supply this quantitative rate.

Fix $c=\lambda^2>1$, $L=\log c$, and $\alpha_j=2\pi j/L$.
Assume the exact real self-adjoint realization on $\ell^2(\mathbb Z)$ is
$A=D_L+K$, with $D_L$ diagonal, even, increasing with $|j|$,
$d_j\sim\log|j|$, $\|K\|\le K_*$, and
$\operatorname{Dom}A=\operatorname{Dom}D_L$.
Assume reflection commutes with $A$, and the off-diagonal entries satisfy
\[
 (\alpha_i-\alpha_j)A_{ij}=g_i-g_j,\qquad g_j=\alpha_j A_{j0}.
\]
Let $a$ be a real even unit ground state \emph{in the even sector}, with
eigenvalue $\ell$, simple and isolated there by a gap $g>0$.
Write $B=A-\ell I$. Thus $B$ is nonnegative on even vectors; no positivity
of its odd restriction is assumed. Let $v_N$ be a unit even ground Ritz
vector on $|j|\le N$, phase aligned with $a$, and set
\[
 \Delta_N=\ell_N-\ell,\qquad Q_N=I-P_N,\qquad
 r_N=(A-\ell_N)v_N=Q_NKv_N.
\]
Fix a positive simple nonlattice real root $\nu=\nu_{c,\infty,k}$ of
$F_a$, and define
\begin{align*}
 b_j&=a_j/(\alpha_j-\nu),&
 b_j^+&=\nu a_j/(\alpha_j^2-\nu^2),\\
 S&=\langle a,b\rangle=\langle a,b^+\rangle,&
 G&=\sum_j g_j b_j,\qquad p=2L^{-1/2}\sin(\nu L/2).
\end{align*}

\begin{proposition}[Exact even adjoint and its positive denominator]
\label{prop:mode-transfer-adjoint}
Under these hypotheses $b\in\operatorname{Dom}A$, $G\ne0$, and
\begin{equation}\label{eq:mode-transfer-exact}
 F_{v_N}(\nu)=\frac pG\left\{
   \Delta_N\langle b^+,v_N\rangle+\langle Q_Nb^+,r_N\rangle\right\}.
\end{equation}
Moreover, with $f=F_a'(\nu)$,
\begin{equation}\label{eq:mode-transfer-positive}
 \langle b,Bb\rangle=\frac{Gf}{p}
     =2\langle b^+,Bb^+\rangle>0.
\end{equation}
In particular the coefficient
\begin{equation}\label{eq:mode-transfer-coefficient}
 \mathcal C_{c,k}=-\frac{pS}{Gf}
       =-\frac{S}{\langle b,Bb\rangle}
\end{equation}
is well defined. No nonvanishing assumption on its numerator is made.
\end{proposition}
\begin{proof}
The multiplier $(\alpha_j-\nu)^{-1}$ is bounded and commutes with $D_L$,
so it preserves its operator domain. The series $\sum b_j$ is absolutely
convergent by Cauchy--Schwarz. So is $G$, since $(A_{j0})_j\in\ell^2$
and $\alpha_j/(\alpha_j-\nu)$ is bounded. Neither $g$ nor the all-ones
sequence is being treated as an $\ell^2$ vector.
The root condition is $\sum_j b_j=0$. Subtracting the $i$th eigenvector
equation after division by $\alpha_i-\nu$, and using the off-diagonal
identity, gives the convergent component identity
\[
 (Bb)_i=\sum_j A_{ij}a_j
 \left(\frac1{\alpha_j-\nu}-\frac1{\alpha_i-\nu}\right)
 =-\frac G{\alpha_i-\nu}.
\]
If $\phi_j=p/(\nu-\alpha_j)$ is the evaluation representer, this says
$Bb=(G/p)\phi$. Taking the even part and pairing with $v_N$ proves
\eqref{eq:mode-transfer-exact} whenever $G\ne0$.

For completeness, nonvanishing and positivity follow in the even sector
itself. Differentiating the secular sum at its root yields
$f/p=-\sum_j a_j/(\alpha_j-\nu)^2$. Evenness and the root condition give
$\sum_j a_j/(\alpha_j^2-\nu^2)=0$. Subtracting the two energies then gives
\[
 \langle b,Bb\rangle-2\langle b^+,Bb^+\rangle
 =-G\sum_j\frac{a_j}{\alpha_j^2-\nu^2}=0.
\]
The even vector $b^+$ is not proportional to $a$: proportionality would
confine the support of $a$ to one value of $|j|$, whose even transform
has no positive nonlattice zero in its secular factor. The simple
even-ground hypothesis therefore gives $\langle b^+,Bb^+\rangle>0$.
The component identity now proves $G\ne0$ and
\eqref{eq:mode-transfer-positive}. This positivity argument concerns
this divided vector; it does not assert positivity of the full operator.
\end{proof}

\begin{proposition}[Projection energy and the omitted adjoint term]
\label{prop:mode-transfer-tail}
Put $t_N=Q_Na$, $\tau_N^2=\|t_N\|^2$, $q_N=\langle t_N,Bt_N\rangle$,
\[
 m_N=d_{N+1}-\ell-K_*,\qquad \chi_N=K_*^2/(g m_N).
\]
If $\tau_N^2<1$, $m_N>0$ and $\chi_N<1$, then
\begin{equation}\label{eq:mode-transfer-energy}
 q_N(1-\chi_N)\le\Delta_N\le q_N/(1-\tau_N^2).
\end{equation}
If also $\alpha_{N+1}>\nu$, then
\begin{align}
 |\langle Q_Nb^+,r_N\rangle|&\le\Delta_N\beta_N,
 \label{eq:mode-transfer-tail}\\
 \beta_N&=\frac{\nu}{\alpha_{N+1}^2-\nu^2}
 \sqrt{\frac{1+2K_*/m_N}{1-K_*^2/(g m_N)}}.
 \nonumber
\end{align}
At fixed $c,k$, $\beta_N=O_{c,k}(N^{-2})$.
\end{proposition}
\begin{proof}
Let $w=P_Na$. Since $Ba=0$, one has $\langle w,Bw\rangle=q_N$
and $P_NBP_Nw=-P_NKt_N$. Expand $w$ in an orthonormal eigenbasis of
$P_NBP_N$: its first eigenvalue is $\Delta_N$, and every other one is
at least $g$ by min--max. Consequently
\[
 q_N\le\Delta_N+\|P_NKt_N\|^2/g
       \le\Delta_N+K_*^2\tau_N^2/g.
\]
On the omitted subspace $B\ge m_N I$, so $\tau_N^2\le q_N/m_N$.
This proves the lower bound; testing $w/\|w\|$ gives the upper bound.

On the omitted subspace write $Q_Nb^+=Mt_N$, where $M$ is the diagonal
multiplier $\nu/(\alpha_j^2-\nu^2)$ and
$\|M\|\le\mu_N:=\nu/(\alpha_{N+1}^2-\nu^2)$.
Since $D_L-\ell I$ is positive there and commutes with $M$,
\[
 \langle Mt_N,BMt_N\rangle
 \le\mu_N^2(q_N+2K_*\tau_N^2)
 \le\mu_N^2\Delta_N\frac{1+2K_*/m_N}{1-\chi_N}.
\]
Both vectors are even. Cauchy--Schwarz in the nonnegative even form $B$
therefore gives
\[
 |\langle Q_Nb^+,Bv_N\rangle|^2
 \le\langle Q_Nb^+,BQ_Nb^+\rangle\Delta_N.
\]
The $\Delta_N v_N$ term vanishes in this pairing, proving the assertion.
\end{proof}

\paragraph{An explicit root remainder.}
Choose $I=[\nu-r,\nu+r]$ on which
$h=\inf_I|F_a'|>0$, and put
$M_1=L^{3/2}/\sqrt{12}$, $M_2=L^{5/2}/\sqrt{80}$.
These bound the first and second Fourier-evaluation derivatives of unit
vectors on the real axis, by the second and fourth moments of the centered
interval. The gap gives
$e_N:=\sqrt{2\Delta_N/g}\ge\|v_N-a\|$. Define
\[
 m=h-M_1e_N,\qquad
 X_N=|p/G|\Delta_N\bigl(|S|+\|b^+\|e_N+\beta_N\bigr).
\]
If $m>0$ and $X_N<mr$, $F_{v_N}$ has exactly one root $\nu_N$ in $I$,
with $|\nu_N-\nu|\le X_N/m$. Taylor's theorem gives
\begin{align}
 |\nu_N-\nu-\mathcal C_{c,k}\Delta_N|
 &\le \frac{|p/G|\Delta_N(\|b^+\|e_N+\beta_N)}{|f|}
 \nonumber\\
 &\quad+\frac{M_1e_NX_N}{m|f|}
       +\frac{M_2X_N^2}{2m^2|f|}.
 \label{eq:mode-transfer-remainder}
\end{align}
Indeed, the derivative of $F_{v_N}$ has fixed sign and magnitude at least
$m$ on $I$, while $|F_{v_N}(\nu)|\le X_N$. The endpoint signs prove
existence and uniqueness. Expand at $\nu$, use
$|F_{v_N}'(\nu)-f|\le M_1e_N$, and apply
\eqref{eq:mode-transfer-exact} to obtain the displayed remainder.
Thus, as $N\to\infty$ with $c,k$ fixed,
\begin{equation}\label{eq:mode-transfer-law}
 \boxed{\nu_{c,N,k}-\nu_{c,\infty,k}
 =\mathcal C_{c,k}\Delta_N
  +O_{c,k}(\Delta_NN^{-2}+\Delta_N^{3/2}).}
\end{equation}
If $S=0$, this formula does not imply a nonzero leading term or equality
of the two convergence rates. The onset $m_N>K_*^2/g$, pole separation,
and the root-derivative margin cannot be suppressed in growing-$c,k$ use.
For example, $A=3I-\mathbf1\mathbf1^T$ on frequencies $-1,0,1$ has a
uniform simple ground state. Its secular root $\nu=1/\sqrt3$ has $S=0$.
This is an exact Loewner example, not a CCM state.

\paragraph{Residuals and arithmetic conditioning.}
For an approximate parent let $r_a=(A-\ell)a$ and
$s=\sum_j a_j/(\alpha_j-\nu)$. The component identity becomes
\[
 Bb=-G(\alpha-\nu)^{-1}+e_b,\qquad
 (e_b)_j=\frac{(r_a)_j+s g_j}{\alpha_j-\nu}.
\]
For any even child, with $r_v=(A-\ell_N)v$, its exact finite algebra gives
\[
 F_v(\nu)=\frac pG\left\{(\ell_N-\ell)\langle b,v\rangle
              +\langle b,r_v\rangle-\langle e_b,v\rangle\right\}.
\]
Using only the omitted block of $r_v$ requires an additional bound on its
retained block. For two scalar terms $Y=q+t\ne0$, changing only $q$ by a
relative $\eta$ changes $Y$ by exactly $\eta q$; the amplification
$|q/(q+t)|$ can be large. It is one component of the precision budget,
not a complete prescription for $P$ or evidence that an input artifact
is corrupt.

\paragraph{Resolution and the remaining parameter law.}
For certified upper bounds on $\Delta_N$ and the explicit remainder,
\[
 B_{\mathrm{mode},k}
 =\frac{|\mathcal C_{c,k}|\overline\Delta_N+
              \overline{\mathcal R}_{c,k,N}}{\gamma_k}
\]
feeds the three-channel sufficient budget in
Section~\ref{ssec:indexed-resolution-budget}. This still depends on the
continuum ground state, its tail, gap and root conditioning; no closed
uniform function of $c,N,P,k$ has been established. In particular
$\Delta_N$ is not the finite eigenvalue $\ell_N$, and $\nu_{c,\infty,k}$
is not being identified with the Riemann ordinate $\gamma_k$.

The intrinsic first attainment at depth $D_*$ is the minimum of
$\{N\in\mathcal A_{c,k}:|\nu_{c,N,k}-\gamma_k|/\gamma_k\le10^{-D_*}\}$,
with value $+\infty$ for an empty set. A persistent threshold instead
requires every later integer to be admissible and accurate. They may
differ under cancellation. The monotone zero-shift family in
Section~\ref{ssec:certified-accuracy-budget} supplies its own special
ordering theorem; that theorem is not automatically transferred to the
changing-ground family considered here.

The admissibility set must specify the root observable. For positive
nonlattice $T$, the finite transform has exactly
\[
 J_{c,N}(T)=\max\{0,\lfloor T L/(2\pi)\rfloor-N\}
\]
forced positive lattice locations below $T$ outside the retained band.
This counts locations, not multiplicity or all transform zeros. A
reference-matched movable root need not have the same ordinal in the
full transform; other roots and separation must also be counted.

\paragraph{Verification scope.}
The displayed arguments were independently reconstructed and checked with
3,900 exact rational identities and 414 separate 80-digit Ritz checks
on 18 finite Loewner fixtures. These fixtures test the algebra, including
the parity energy identity and retained-residual defects. They are not
CCM assemblies and do not prove a uniform arithmetic tail estimate.
Finite-parent numerical comparisons do not replace the continuum
hypotheses or the explicit omitted-mode term in this theorem.

\paragraph{Reproduced finite-parent comparisons.}
A separate retained-data replay uses the stored $c=1000$, $N=800,890$
states and the same 1700-digit analytic generator. The exact two-term
identity agrees at all 13 fixed indices to relative discrepancy below
$2\cdot10^{-397}$. Nine pairs meet the frozen accuracy and root-residual
gate. The full two-term Newton displacement passes all nine one-percent
tests; deleting the omitted-mode term or using the parent coefficient
alone passes none. At $k=1$ omission already causes $-1.60335\%$ error.
This is consistent with a fixed-cutoff late theorem whose onset is not
yet attained in the large-cutoff ramp.

The measured relative depths make the indexed resolution dependence clear:
\begin{center}\small
\begin{tabular}{@{}rrr@{}}
\toprule
Target label $k$ & $N=800$ & $N=890$\\
\midrule
1 & 1260.364036 & 1358.789387\\
100 & 1088.105848 & 1190.938797\\
300 & 729.949771 & 862.073328\\
500 & 73.905037 & 299.440729\\
525 & 29.817503 & 180.934808\\
\bottomrule
\end{tabular}
\end{center}
These are reference-assisted point matches, not certified full-spectrum
ordinals or minima over untested integers. At the parent target heights
labeled $500,600,610$, respectively $1,142,155$ forced lattice locations
already lie below them outside the retained band. The labels therefore
require their explicit movable-root convention.

Four late comparisons at $c=5,9$, using $320\to640$ and $640\to1280$,
instead pass the one-percent parent-coefficient test. Their overlapping
states and 45-significant-digit coefficient summaries are retained with
summary-rounding allowances. The original interval-center eigenvalue
calculation, which fails the high-index identity, is also retained;
Rayleigh products on the same vectors and generator repair it. Replaying
these computations, including both versions, reproduces all eight
comparison tables byte for byte and the selected exact-decimal identity
fields. No source assembly or reference-zero certification is inferred
from this replay. A separate contraction audit has two failures among
11 retained late rows, so a universal one-half gain-ratio shortcut is
not supported by that corpus.
